% Options for packages loaded elsewhere
\PassOptionsToPackage{unicode}{hyperref}
\PassOptionsToPackage{hyphens}{url}
\documentclass[
  ignorenonframetext,
]{beamer}
\newif\ifbibliography
\usepackage{pgfpages}
\setbeamertemplate{caption}[numbered]
\setbeamertemplate{caption label separator}{: }
\setbeamercolor{caption name}{fg=normal text.fg}
\beamertemplatenavigationsymbolsempty
% remove section numbering
\setbeamertemplate{part page}{
  \centering
  \begin{beamercolorbox}[sep=16pt,center]{part title}
    \usebeamerfont{part title}\insertpart\par
  \end{beamercolorbox}
}
\setbeamertemplate{section page}{
  \centering
  \begin{beamercolorbox}[sep=12pt,center]{section title}
    \usebeamerfont{section title}\insertsection\par
  \end{beamercolorbox}
}
\setbeamertemplate{subsection page}{
  \centering
  \begin{beamercolorbox}[sep=8pt,center]{subsection title}
    \usebeamerfont{subsection title}\insertsubsection\par
  \end{beamercolorbox}
}
% Prevent slide breaks in the middle of a paragraph
\widowpenalties 1 10000
\raggedbottom
\AtBeginPart{
  \frame{\partpage}
}
\AtBeginSection{
  \ifbibliography
  \else
    \frame{\sectionpage}
  \fi
}
\AtBeginSubsection{
  \frame{\subsectionpage}
}
\usepackage{iftex}
\ifPDFTeX
  \usepackage[T1]{fontenc}
  \usepackage[utf8]{inputenc}
  \usepackage{textcomp} % provide euro and other symbols
\else % if luatex or xetex
  \usepackage{unicode-math} % this also loads fontspec
  \defaultfontfeatures{Scale=MatchLowercase}
  \defaultfontfeatures[\rmfamily]{Ligatures=TeX,Scale=1}
\fi
\usepackage{lmodern}
\ifPDFTeX\else
  % xetex/luatex font selection
\fi
% Use upquote if available, for straight quotes in verbatim environments
\IfFileExists{upquote.sty}{\usepackage{upquote}}{}
\IfFileExists{microtype.sty}{% use microtype if available
  \usepackage[]{microtype}
  \UseMicrotypeSet[protrusion]{basicmath} % disable protrusion for tt fonts
}{}
\makeatletter
\@ifundefined{KOMAClassName}{% if non-KOMA class
  \IfFileExists{parskip.sty}{%
    \usepackage{parskip}
  }{% else
    \setlength{\parindent}{0pt}
    \setlength{\parskip}{6pt plus 2pt minus 1pt}}
}{% if KOMA class
  \KOMAoptions{parskip=half}}
\makeatother
\usepackage{color}
\usepackage{fancyvrb}
\newcommand{\VerbBar}{|}
\newcommand{\VERB}{\Verb[commandchars=\\\{\}]}
\DefineVerbatimEnvironment{Highlighting}{Verbatim}{commandchars=\\\{\}}
% Add ',fontsize=\small' for more characters per line
\newenvironment{Shaded}{}{}
\newcommand{\AlertTok}[1]{\textcolor[rgb]{1.00,0.00,0.00}{\textbf{#1}}}
\newcommand{\AnnotationTok}[1]{\textcolor[rgb]{0.38,0.63,0.69}{\textbf{\textit{#1}}}}
\newcommand{\AttributeTok}[1]{\textcolor[rgb]{0.49,0.56,0.16}{#1}}
\newcommand{\BaseNTok}[1]{\textcolor[rgb]{0.25,0.63,0.44}{#1}}
\newcommand{\BuiltInTok}[1]{\textcolor[rgb]{0.00,0.50,0.00}{#1}}
\newcommand{\CharTok}[1]{\textcolor[rgb]{0.25,0.44,0.63}{#1}}
\newcommand{\CommentTok}[1]{\textcolor[rgb]{0.38,0.63,0.69}{\textit{#1}}}
\newcommand{\CommentVarTok}[1]{\textcolor[rgb]{0.38,0.63,0.69}{\textbf{\textit{#1}}}}
\newcommand{\ConstantTok}[1]{\textcolor[rgb]{0.53,0.00,0.00}{#1}}
\newcommand{\ControlFlowTok}[1]{\textcolor[rgb]{0.00,0.44,0.13}{\textbf{#1}}}
\newcommand{\DataTypeTok}[1]{\textcolor[rgb]{0.56,0.13,0.00}{#1}}
\newcommand{\DecValTok}[1]{\textcolor[rgb]{0.25,0.63,0.44}{#1}}
\newcommand{\DocumentationTok}[1]{\textcolor[rgb]{0.73,0.13,0.13}{\textit{#1}}}
\newcommand{\ErrorTok}[1]{\textcolor[rgb]{1.00,0.00,0.00}{\textbf{#1}}}
\newcommand{\ExtensionTok}[1]{#1}
\newcommand{\FloatTok}[1]{\textcolor[rgb]{0.25,0.63,0.44}{#1}}
\newcommand{\FunctionTok}[1]{\textcolor[rgb]{0.02,0.16,0.49}{#1}}
\newcommand{\ImportTok}[1]{\textcolor[rgb]{0.00,0.50,0.00}{\textbf{#1}}}
\newcommand{\InformationTok}[1]{\textcolor[rgb]{0.38,0.63,0.69}{\textbf{\textit{#1}}}}
\newcommand{\KeywordTok}[1]{\textcolor[rgb]{0.00,0.44,0.13}{\textbf{#1}}}
\newcommand{\NormalTok}[1]{#1}
\newcommand{\OperatorTok}[1]{\textcolor[rgb]{0.40,0.40,0.40}{#1}}
\newcommand{\OtherTok}[1]{\textcolor[rgb]{0.00,0.44,0.13}{#1}}
\newcommand{\PreprocessorTok}[1]{\textcolor[rgb]{0.74,0.48,0.00}{#1}}
\newcommand{\RegionMarkerTok}[1]{#1}
\newcommand{\SpecialCharTok}[1]{\textcolor[rgb]{0.25,0.44,0.63}{#1}}
\newcommand{\SpecialStringTok}[1]{\textcolor[rgb]{0.73,0.40,0.53}{#1}}
\newcommand{\StringTok}[1]{\textcolor[rgb]{0.25,0.44,0.63}{#1}}
\newcommand{\VariableTok}[1]{\textcolor[rgb]{0.10,0.09,0.49}{#1}}
\newcommand{\VerbatimStringTok}[1]{\textcolor[rgb]{0.25,0.44,0.63}{#1}}
\newcommand{\WarningTok}[1]{\textcolor[rgb]{0.38,0.63,0.69}{\textbf{\textit{#1}}}}
\usepackage{longtable,booktabs,array}
\newcounter{none} % for unnumbered tables
\usepackage{calc} % for calculating minipage widths
\usepackage{caption}
% Make caption package work with longtable
\makeatletter
\def\fnum@table{\tablename~\thetable}
\makeatother
\setlength{\emergencystretch}{3em} % prevent overfull lines
\providecommand{\tightlist}{%
  \setlength{\itemsep}{0pt}\setlength{\parskip}{0pt}}
% Slide layout defaults for warning-clean scientific decks.
\usepackage{etoolbox}
\IfFileExists{xurl.sty}{\usepackage{xurl}}{}
\IfFileExists{seqsplit.sty}{\usepackage{seqsplit}}{\newcommand{\seqsplit}[1]{#1}}
\protected\def\breakseq#1{\seqsplit{#1}}
\protected\def\breaktt#1{\begingroup\ttfamily\seqsplit{#1}\endgroup}
\setlength{\emergencystretch}{6em}
\tolerance=5000
\hbadness=10000
\hfuzz=1pt
\setlength{\tabcolsep}{2pt}
\AtBeginEnvironment{longtable}{\tiny\renewcommand{\arraystretch}{0.86}\setlength{\tabcolsep}{1pt}}
\AtBeginEnvironment{tabular}{\tiny\renewcommand{\arraystretch}{0.86}\setlength{\tabcolsep}{1pt}}

% Natbib and cross-reference fallbacks — slides don't load natbib
% or cleveref, but combined-PDF manuscript prose may emit these
% commands. The fallback renders citations as a bracketed key list
% and cross-references as detokenized labels so slides don't fail on
% undefined control sequences or raw underscores. \providecommand is
% a no-op if packages load later.
\providecommand{\citep}[1]{[#1]}
\providecommand{\citet}[1]{#1}
\providecommand{\citealp}[1]{#1}
\providecommand{\citeauthor}[1]{#1}
\providecommand{\citeyear}[1]{#1}
\providecommand{\cref}[1]{\texttt{\detokenize{#1}}}
\providecommand{\Cref}[1]{\texttt{\detokenize{#1}}}
\usepackage{bookmark}
\IfFileExists{xurl.sty}{\usepackage{xurl}}{} % add URL line breaks if available
\urlstyle{same}
% fallback for those not using the hyperref driver hyperxmp:
\makeatletter
\@ifundefined{xmpquote}{\newcommand{\xmpquote}[1]{#1}}{}
\makeatother
\hypersetup{
  hidelinks,
  pdfcreator={LaTeX via pandoc}}

\author{\texorpdfstring{}{}}
\date{}

\begin{document}

\section{Methods}\label{methods}

\begin{frame}[fragile,allowframebreaks]{Methods}
The \texttt{template/} architecture is deliberately bifurcated into a
globally shared \texttt{infrastructure/} layer and project-specific
\texttt{projects/} silos. This section describes the four core design
patterns, the YAML-declared pipeline (14 stages; default full 10)
pipeline that operationalizes them, and the AI collaboration model that
distinguishes this system from conventional research templates.
\end{frame}

\begin{frame}[fragile,allowframebreaks]{The Two-Layer Architecture}
\protect\phantomsection\label{the-two-layer-architecture}
The repository is organized into two strictly separated layers:

\textbf{Infrastructure Layer} (\texttt{infrastructure/}): 23
infrastructure subdirectories---20 of them independently-importable
Python packages---comprising \textasciitilde616 modules and providing
reusable services. Each importable package has its own
\texttt{\_\_init\_\_.py}, \texttt{AGENTS.md}, and \texttt{README.md},
and exports a well-defined public API (the remaining subdirectories,
e.g.~\texttt{config/}, hold shared configuration). The infrastructure
layer knows nothing about any specific project---it provides generic
capabilities (logging, rendering, validation, steganography) that any
project may consume.

\textbf{Project Layer} (\texttt{projects/}): Self-contained research
workspaces. Each project directory contains:

{\def\LTcaptype{none} % do not increment counter
\begin{longtable}[]{@{}ll@{}}
\toprule\noalign{}
Directory & Purpose \\
\midrule\noalign{}
\endhead
\texttt{manuscript/} & Markdown chapters and \texttt{config.yaml} \\
\texttt{scripts/} & Thin orchestrator scripts (Stage 02) \\
\texttt{src/} & Project-specific Python modules \\
\texttt{tests/} & Project-specific test suite \\
\texttt{data/} & Input datasets and generated data \\
\texttt{output/} & Pipeline artifacts: PDF, figures, reports, logs \\
\texttt{docs/} & Project-specific architecture documentation \\
\bottomrule\noalign{}
\end{longtable}
}

The two layers communicate exclusively through Python imports and
filesystem paths. No project modifies infrastructure code; no
infrastructure module references a specific project by name (except via
runtime project discovery).
\end{frame}

\begin{frame}[fragile,allowframebreaks]{The Standalone Project Paradigm}
\protect\phantomsection\label{the-standalone-project-paradigm}
Projects are designed to be completely self-contained. Adding a new
project requires no changes to the infrastructure layer, no
modifications to \texttt{pyproject.toml}, and no updates to the pipeline
orchestrator. A project is automatically discovered if and only if it
satisfies two conditions:

\begin{enumerate}
\tightlist
\item
  It exists as a subdirectory of \texttt{projects/}.
\item
  It contains the file \breaktt{manuscript/config.yaml}.
\end{enumerate}

This paradigm enables horizontal scaling: N researchers can maintain N
independent projects within a single repository, sharing infrastructure
without coupling. Each project declares its own testing tolerances,
manuscript metadata, LLM review preferences, and rendering configuration
in its \texttt{config.yaml}. The system currently hosts its public
canonical exemplars under \breaktt{projects/templates/}
(\breaktt{templates/template\_active\_inference},
\breaktt{templates/template\_autoresearch\_project},
\breaktt{templates/template\_autoscientists},
\breaktt{templates/template\_code\_project},
\breaktt{templates/template\_eda\_notebook},
\breaktt{templates/template\_gold\_refinement},
\breaktt{templates/template\_literature\_meta\_analysis},
\breaktt{templates/template\_madlib},
\breaktt{templates/template\_methods\_paper},
\breaktt{templates/template\_newspaper},
\breaktt{templates/template\_prose\_project},
\breaktt{templates/template\_search\_project},
\breaktt{templates/template\_sia},
\breaktt{templates/template\_storybook},
\breaktt{templates/template\_template},
\breaktt{templates/template\_textbook}), including this meta-manuscript
at \breaktt{projects/templates/template\_template/}.
\end{frame}

\begin{frame}[fragile,allowframebreaks]{The Thin Orchestrator Pattern}
\protect\phantomsection\label{the-thin-orchestrator-pattern}
All scripts in \texttt{scripts/} (both infrastructure-level and
project-level) follow the Thin Orchestrator pattern
\breakseq{[@gamma1995design]}:

\begin{itemize}
\tightlist
\item
  \textbf{No domain logic}: Scripts contain zero algorithmic
  implementation. They import functions from \texttt{src/} modules and
  wire them to infrastructure services.
\item
  \textbf{Configuration-driven}: Behavior is parameterized by
  \texttt{config.yaml}, not by hardcoded values.
\item
  \textbf{Stateless}: Scripts read inputs, call functions, write
  outputs. They maintain no persistent state between invocations.
\item
  \textbf{Logged}: Every significant action is logged via
  \breaktt{infrastructure.core.logging.utils.get\_logger}.
\end{itemize}

This pattern ensures that all testable logic lives in \texttt{src/}
where it is subject to the Zero-Mock testing policy, while scripts
remain thin enough to be audited by visual inspection. The separation
draws on the Mediator pattern from Gamma et al. \breakseq{[@gamma1995design]},
where scripts mediate between infrastructure services and
project-specific code without implementing any logic of their own.

To make this concrete, the following contrasts the anti-pattern with the
correct pattern:

\begin{Shaded}
\begin{Highlighting}[]
\CommentTok{\# ANTI{-}PATTERN: domain logic embedded in script}
\KeywordTok{def}\NormalTok{ calculate\_average(data):          }\CommentTok{\# ← never put computation here}
    \ControlFlowTok{return} \BuiltInTok{sum}\NormalTok{(data) }\OperatorTok{/} \BuiltInTok{len}\NormalTok{(data)}

\NormalTok{result }\OperatorTok{=}\NormalTok{ calculate\_average([}\DecValTok{1}\NormalTok{, }\DecValTok{2}\NormalTok{, }\DecValTok{3}\NormalTok{])}

\CommentTok{\# CORRECT PATTERN: script imports from src/ and only wires}
\ImportTok{from}\NormalTok{ projects.my\_project.src.statistics }\ImportTok{import}\NormalTok{ calculate\_average}

\NormalTok{result }\OperatorTok{=}\NormalTok{ calculate\_average([}\DecValTok{1}\NormalTok{, }\DecValTok{2}\NormalTok{, }\DecValTok{3}\NormalTok{])  }\CommentTok{\# ← scripts wire, never compute}
\end{Highlighting}
\end{Shaded}

The critical property is that \breaktt{calculate\_average} in the correct
pattern lives in a testable \texttt{src/} module, is covered by the
Zero-Mock test suite, and can be independently imported, tested, and
reused---whereas the anti-pattern buries logic in a script that is
invisible to coverage tools.
\end{frame}

\end{document}
